%% generated by tools/make_tables.py -- do not edit
\begin{table}[t]
\centering
\caption{观测者距离扫描：在运行的应用里把相机半径从 $8M$ 推到 $150M$，每一档都用内建探针反解阴影半径并与式~\eqref{eq:shadowpx} 的闭式比较（窗口 $1600\times900$、$\tan\psi$ 视场 $58^\circ$、$\mathtt{autoQuality}$ 关闭、单帧冻结）。闭式在几何单位下为 $R_{\rm px}=f_{\rm px}\tan\psi_{\rm sh}$，只依赖 $r_0/M$；实测相对偏差在 $\pm0.13\%$ 以内，绝对残差 $\le0.11$\,px。覆盖三个数量级的 $r_0$ 上残差都保持在百分之几个像素，说明限制这一测量的是覆盖率量化带 $\sigma_{\rm px}=0.0147$\,px 与掩膜本身亚像素级的系统项，而不是阴影角半径的闭式。数据由 \texttt{tools/verify.py} 采集（GPU：ANGLE (Intel, Intel(R) UHD Graphics (0x000046D2) Direct3D11 vs\_5\_0 ps\_5\_0, D3D11)）。}
\label{tab:r0scan}
\small
\begin{tabular}{rrrrr}
\toprule
$r_0/M$ & 实测 [px] & 闭式 [px] & 相对 [\%] & 残差 [px] \\
\midrule
8 & 552.33 & 552.31 & +0.004 & +0.02 \\
12 & 349.38 & 349.35 & +0.006 & +0.02 \\
20 & 206.46 & 206.46 & -0.002 & -0.01 \\
26 & 158.85 & 158.83 & +0.012 & +0.02 \\
40 & 103.52 & 103.62 & -0.098 & -0.10 \\
80 & 52.17 & 52.17 & -0.013 & -0.01 \\
150 & 27.92 & 27.95 & -0.122 & -0.03 \\
\bottomrule
\end{tabular}

\end{table}
